\chapter{Gaussian and Eisenstein depth}\label{ch:04-depth}
The classical transformations suggest a sharper question: how much nesting
is necessary after known identities have been used? We first prove an
infinite family of reductions. The other families illustrate how a successful
experimental identity gives an upper bound on depth, while a failed search
leaves open both a large-coefficient reduction and a missing argument family.

\section{Depth conventions and Gaussian generators}
Write $G=\beta(2)$, $\beta(s)=\sum_{n\ge0}(-1)^n/(2n+1)^s$, and
\[
g_{a,b}=\Im\Li_{a,b}(i,1),\qquad h_{a,b}=\Re\Li_{a,b}(i,1),\qquad
\lambda_n=\Im\Li_n((1+i)/2).
\]
Subscripts such as $g_{41}$ mean $g_{4,1}$.
Series depth counts nested indices. In the experimental tables, a
``depth-one reduction'' means a polynomial in the stated single-value
basket, which is a different convention from the additive depth filtration
on a motivic algebra. Products of two depth-two factors have additive depth
at most four; their membership in an algebra generated by doubles does not
make them depth at most two in that additive filtration.

\section{An infinite even-weight reduction}
Define
\[
S_p=\sum_{n\ge0}\frac{(-1)^nH_n}{(2n+1)^p},\qquad p\ge1.
\]
For $p=1$ the sum is conditionally convergent; for $p>1$ it is absolute.
\begin{theorem}[Odd-denominator harmonic sums at even weight]
For every integer $m\ge0$,
\begin{equation}\label{gauss:eq:even-general}
S_{2m+1}=(2m+1)\beta(2m+2)
-[t^{2m}]\frac{\pi}{4\cos(\pi t/2)}
\left(2\log2+\sum_{j\ge1}(2-2^{-2j})\zeta(2j+1)t^{2j}\right).
\end{equation}
In particular, each sum of even weight $2m+2$ belongs to the algebra
generated by $\pi,\log2$, odd zeta values and even beta values.
\end{theorem}
\begin{proof}
Put
\[
I(t)=\int_0^1\frac{x^t\log(1+x^2)}{1+x^2}\,dx,\quad
J(t)=\int_0^1\frac{x^t}{1+x^2}\,dx,\quad
A(t)=\int_0^\infty\frac{x^t\log(1+x^2)}{1+x^2}\,dx.
\]
These are analytic near $t=0$. Splitting $A$ at one and substituting
$x=1/u$ in the second part gives
$A(t)=I(t)+I(-t)-2J'(-t)$.
The generating function $\sum H_nz^n=-\log(1-z)/(1-z)$,
used first with an Abel factor and then by dominated convergence, gives
$I^{(2m)}(0)=-(2m)!S_{2m+1}$.
Also $J^{(2m+1)}(0)=-(2m+1)!\beta(2m+2)$, by the geometric series for
$1/(1+x^2)$; its integrated absolute majorant is summable.
Consequently
\[
S_{2m+1}=(2m+1)\beta(2m+2)-\frac{[t^{2m}]A(t)}2.
\]
Differentiating the beta integral in its second parameter yields
\[
A(t)=\frac{\pi}{2\cos(\pi t/2)}
\left[\psi(1)-\psi\left(\frac{1-t}{2}\right)\right],\qquad |\Re t|<1.
\]
Its even part follows from $\psi(1)-\psi(1/2)=2\log2$ and
$\psi^{(2j)}(1/2)=-(2j)!(2^{2j+1}-1)\zeta(2j+1)$.
Extracting the coefficient proves the formula.
\end{proof}
The first cases are
\begin{align}
S_1&=G-\frac\pi2\log2,\\
S_3&=3\beta(4)-\frac{\pi^3}{16}\log2-\frac{7\pi}{16}\zeta(3),\\
768S_5&=3840\beta(6)-5\pi^5\log2-42\pi^3\zeta(3)-372\pi\zeta(5),\\
92160S_7&=645120\beta(8)-61\pi^7\log2-525\pi^5\zeta(3)
-5580\pi^3\zeta(5)-45720\pi\zeta(7).
\end{align}
The source's high-precision experiments suggested the last two identities;
the generating integral proves them and every higher odd-index case at once.
The proof gives a sufficient reduction. It does not prove that the
even beta values are independent, nor that the odd-weight sums have minimal
depth two. No such conclusion follows merely from the absence of a new
odd beta value.

\section{Candidate Gaussian reductions at weights five and six}
The following equations collect the source's accepted numerical relations.
They include the corrected dependence among the four weight-five generators,
the five weight-six double reductions, and the triple reductions through
weight five. Historical residuals range from about $10^{-193}$ to
$10^{-690}$. The source scripts for every such relation are not present in
this checkout, so these remain numerically supported identities unless an
analytic derivation is explicitly supplied. In particular the weight-five
relation gives a spanning set of at most three directions, rather than an
independence theorem.
\begin{equation}
  \label{gauss:eq:S4-closed}
  S_4 \;=\; \sum_{n\ge0}\frac{(-1)^n H_n}{(2n+1)^4}
  \;=\; \frac{4\,g_{41}-3\,g_{32}-9\,g_{23}}{7}
        \;+\;\frac{\pi^5}{224}
        \;-\;\frac{27}{224}\,G\,\zeta(3)
        \;-\;2\,\beta(4)\log 2,
\end{equation}

\begin{equation}
  \label{gauss:eq:wt5-sporadic}
  576\,g_{41}+288\,g_{32}+736\,g_{23}+960\,g_{14}
  \;=\; 21\,G\,\zeta(3)-480\,\beta(4)\log 2,
\end{equation}

\begin{align}
  \label{gauss:eq:g51}
  2048\,g_{51} &= -64\pi^3\zeta(3)-527\pi\zeta(5)+4096\,\beta(6),\\
  \label{gauss:eq:g42}
  1536\,g_{42} &= 96\pi^3\zeta(3)-32\pi^2\beta(4)+1581\pi\zeta(5)-8448\,\beta(6),\\
  \label{gauss:eq:g33}
  1024\,g_{33} &= -3\pi^3\zeta(3)+64\pi^2\beta(4)-1581\pi\zeta(5)+4608\,\beta(6),\\
  \label{gauss:eq:g24}
  23040\,g_{24} &= -14\pi^4 G+135\pi^3\zeta(3)-1440\pi^2\beta(4)+23715\pi\zeta(5)-69120\,\beta(6),\\
  \label{gauss:eq:g15}
  92160\,g_{15} &= -150\pi^5\log 2+56\pi^4 G-270\pi^3\zeta(3)+1920\pi^2\beta(4)-675\pi\zeta(5).
\end{align}

\begin{align}
  \Im\Li_{2,1,1}(\iu,1,1)
    &= \lambda_4 + \tfrac12\,\lambda_3\log 2
       - \frac{G\,(\pi^2-4\log^2 2)}{32}
       - \frac{\pi\,(8\log^3 2 + 105\,\zeta(3))}{768},\\[2pt]
  \Im\Li_{1,2,1}(\iu,1,1)
    &= -3\lambda_4 - \lambda_3\log 2
       + \frac{G\,(\pi^2-4\log^2 2)}{32}
       - \frac{3\pi^3\log 2}{256}
       + \frac{\pi\,(2\log^3 2 + 67\,\zeta(3))}{128},\\[2pt]
  \Im\Li_{1,1,2}(\iu,1,1)
    &= 3\lambda_4 + \tfrac12\,\lambda_3\log 2
       + \frac{5\pi^3\log 2}{192}
       - \frac{163\,\pi\,\zeta(3)}{256}.
\end{align}

\section{The colored stuffle identity}
Whenever the defining sums are absolutely convergent, sorting the two
indices in a product proves
\begin{equation}\label{double:eq:stuffle-d1}
\Li_a(x)\Li_b(y)=\Li_{a,b}(x,y)+\Li_{b,a}(y,x)+\Li_{a+b}(xy).
\end{equation}
At conditionally convergent boundary points, Abel limits give the same
identity when the limiting products and constituent series exist.

\section{Gaussian doubles: closed forms}\label{double:sec:gauss-fam}

Write $\Li_{a,b}$ at the three Gaussian base points beyond $(\iu,1)$:
\[
  (\iu,\iu),\qquad (-1,\iu),\qquad (\iu,-1),
\]
all generated to $700$ digits by an accelerated evaluator (Section~\ref{ch:10-discovery}) and agreeing
with Wolfram~$15$ to full precision. Under the level-$4$ complex conjugation $\iu\mapsto-\iu$ the
imaginary parts are odd and the real parts even.

\subsection{Family 1: the stuffle closed forms}

Specializing \eqref{double:eq:stuffle-d1} at the Gaussian points ($uv=-1$ at $(\iu,\iu)$; $uv=-\iu$ at the
mixed pair) gives a complete set of exact closed forms for the symmetric combinations:
\begin{align}
  \Li_{a,b}(\iu,\iu)+\Li_{b,a}(\iu,\iu)   &= \Li_a(\iu)\,\Li_b(\iu)   -\Li_{a+b}(-1),\label{double:eq:g-sym-ii}\\
  \Li_{a,b}(-1,\iu)+\Li_{b,a}(\iu,-1)     &= \Li_a(-1)\,\Li_b(\iu)    -\Li_{a+b}(-\iu),\label{double:eq:g-sym-mix}\\
  \Li_{a,b}(\iu,-1)+\Li_{b,a}(-1,\iu)     &= \Li_a(\iu)\,\Li_b(-1)    -\Li_{a+b}(-\iu),\nonumber
\end{align}
with the clean diagonal corollary
\begin{equation}\label{double:eq:g-diag}
  \Li_{a,a}(\iu,\iu)=\tfrac12\!\left(\Li_a(\iu)^2-\Li_{2a}(-1)\right).
\end{equation}
Every single polylogarithm on the right is elementary in the level-$4$ vocabulary:
$\Imag\Li_n(\iu)=\beta(n)$ is the Dirichlet beta value, $\Real\Li_n(\iu)$ and $\Li_n(-1)=-\eta(n)$
are rational multiples of $\zeta(n)$ (or $\pi^n$ at even $n$), and $\Li_n(-\iu)=\overline{\Li_n(\iu)}$.
So Family~1 reduces every symmetric and diagonal Gaussian double to single $L$-values in closed form.

\begin{remark}
The three stuffle identities \eqref{double:eq:g-sym-ii}--\eqref{double:eq:g-sym-mix} were checked against the stored
$700$-digit values to ${\sim}10^{-692}$ for every admissible $a,b$, and the closed forms were
confirmed symbolically and numerically by Wolfram~$15$ at low weight, including the weight-$6$
diagonal $\Li_{3,3}(\iu,\iu)=\tfrac12(\Li_3(\iu)^2-\Li_6(-1))$, for which the kernel returns a form in
$\beta(4)$, $\pi^4$, $\zeta(3)$.
\end{remark}

The consequence is structural: all depth-$2$ content of the Gaussian doubles lives in the
\emph{antisymmetric} parts $\Li_{a,b}(P)-\Li_{b,a}(P)$, which vanish on the diagonal. These are what
the double shuffle of the next section reduces.

\subsection{Family 2: antisymmetric reductions to the $(\iu,1)$ algebra}

Against the basis of single-polylog monomials together with the $(\iu,1)$ generators
$g_{a,b}=\Imag\Li_{a,b}(\iu,1)$, $h_{a,b}=\Real\Li_{a,b}(\iu,1)$, the \emph{low-weight} antisymmetric
parts reduce, each reduction a new closed-form identity found by integer-relation search and confirmed
in Wolfram~$15$. A representative instance, with $\mu_n=\Real\Li_n(\tfrac{1+\iu}2)$ and
$\lambda_n=\Imag\Li_n(\tfrac{1+\iu}2)$:
\begin{equation}\label{double:eq:fam2-example}
  \Real\!\left[\Li_{1,2}(\iu,\iu)-\Li_{2,1}(\iu,\iu)\right]
  =3\zeta(3)-4\log2\,\mu_2-\pi\lambda_2-4h_{1,1}\log2+3h_{1,2}.
\end{equation}
Such reductions exist at weights~$2$--$4$ at the mixed points and at weight~$3$ at $(\iu,\iu)$; at
higher weight the antisymmetric parts do not reduce \emph{against the $(\iu,1)$ basis alone}. To go
further---and to learn how many are genuinely new---one needs the shuffle.

\section{Gaussian doubles: the double-shuffle dimension}\label{double:sec:gauss-frontier}

\subsection{A cautionary integer-relation count}

A first attempt to count the genuinely-new depth-$2$ directions proceeds by integer-relation search:
test each antisymmetric part for membership in the span of the $(\iu,1)$ generators and single-polylog
products, and tabulate those that fail. Carried out at $700$ digits, this reports \textbf{30}
independent new directions through weight~$6$. The number is wrong. It is an artifact of the spurious
scale: the genuine reductions are of large height, and against a basis of a hundred-odd atoms the
search cannot see them. (Three searches against bases of $46$, $55$, and the $(\iu,1)$-plus-single set
returned ``dimensions'' of $3$, $6$, and $30$ respectively---all different, all wrong, the count
\emph{decreasing} as the basis grew and the ceiling fell.)

\subsection{The shuffle, and the correct count}

The shuffle settles it. The key relation is that the shuffle of a $(-1)$-polylogarithm with an
$\iu$-polylogarithm produces $(\iu,\iu)$ and $\overline{(-1,\iu)}$ doubles, for instance
\begin{equation}\label{double:eq:shuffle-22}
  \Li_2(-1)\,\Li_2(\iu)=\Li_{2,2}(\iu,\iu)+2\Li_{3,1}(\iu,\iu)
     +\overline{\Li_{2,2}(-1,\iu)}+2\,\overline{\Li_{3,1}(-1,\iu)}
\end{equation}
(verified to ${\sim}10^{-694}$), whereas the \emph{stuffle} of the same product produces the
$(-1,\iu),(\iu,-1)$ doubles. Equating ties all three Gaussian points together. Collecting all $240$
stuffle and shuffle relations from the products $\Li_p(s)\Li_q(t)$, $s,t\in\{\iu,-1,-\iu\}$,
$2\le p+q\le6$, verifying each to ${\sim}10^{-600}$, and computing the exact rational rank of the
linear system on the doubles, modulo single polylogarithms and the $(\iu,1)$ generators, gives:

\begin{center}
\renewcommand{\arraystretch}{1.2}
\begin{tabular}{c c c c}
\hline
weight & unknowns (Re$+$Im) & dim (stuffle only) & dim (double shuffle)\\
\hline
$4$ & $18$ & $8$  & $\mathbf 1$\\
$5$ & $24$ & $12$ & $\mathbf 2$\\
$6$ & $30$ & $14$ & $\mathbf 2$\\
\hline
total & & $34$ & $\mathbf 5$\\
\hline
\end{tabular}
\end{center}

\begin{observation}\label{double:thm:gauss-5}
Modulo single polylogarithms and the $(\iu,1)$ generators, the Gaussian double points
$(\iu,\iu),(-1,\iu),(\iu,-1)$ contribute exactly $\mathbf 5$ candidate depth-two generators through
weight~$6$---one at weight~$4$, two at weight~$5$, two at weight~$6$. A representative basis is
\[
  \Real\Li_{3,1}(\iu,-1);\quad \Imag\Li_{3,2}(\iu,-1),\ \Imag\Li_{4,1}(\iu,-1);\quad
  \Real\Li_{4,2}(\iu,-1),\ \Real\Li_{5,1}(\iu,-1).
\]
Every other Gaussian double of weight $\le6$ is an explicit combination of these five, the $(\iu,1)$
generators, and products of single $L$-values.
\end{observation}

The rank is verified two ways: each relation reproduces its value to $10^{-600}$, and the
exact-rational rank is precision-independent. As a final gold-check at weight~$4$, fixing the single
survivor $\Real\Li_{3,1}(\iu,-1)$ to its true value and solving the system predicts all other $17$
weight-$4$ components to ${\sim}10^{-692}$.

\begin{remark}[the methodological moral]
Observation~\ref{double:thm:gauss-5} corrects the integer-relation count of $30$ by a factor of six. The
reductions are exact consequences of the small-height double-shuffle generators, but the
\emph{net} reductions are of large height---so high that an integer-relation search against any
adequate basis cannot find them. This is the recurring trap of high-precision experimental
mathematics, here biting in the direction of \emph{over}-stating independence. The derived double
shuffle, with its exact rank, is immune.
\end{remark}

\section{Triple sums and the scope of formal quotient dimensions}
The word for $\Li_{a,b,c}(i,1,1)$ is
$0^{a-1}(-i)0^{b-1}(-i)0^{c-1}(-i)$ in the outer-letter-first convention.
Shuffle relations express products of lower-weight functions as sums of
such triples. At weight six the source reports rank seven for thirteen
relations on ten triple coordinates, leaving three formal coordinates
represented by $(3,1,2),(3,2,1),(4,1,1)$.
Rank seven means that seven coordinates can be eliminated using the
chosen relations and the three remaining coordinates. It does not mean
that seven particular triples are individually products, or that the
remaining three numerical values are independent.

The source also claimed binomial dimensions $\binom np$ for the actual
depth filtration from a total Hilbert series $1/(1-2t)$.
A total Hilbert series alone cannot determine a depth grading. A grading by
the length of a chosen motivic word basis must first be identified with the
filtration under discussion; moreover a motivic lower bound does not give
a numerical lower bound without injectivity of the period map.
We therefore make no unconditional claim that genuine numerical depth three
first occurs at weight six. A rigorous reduction or obstruction for these
specific triples is an open problem. Panzer's parity theorem supplies exact
lower-depth conjugation combinations, but does not assert independence of
the complementary components.

\section{The Eisenstein analogue}\label{double:sec:eisenstein}

Everything above has a level-$3$ counterpart at the cube roots of unity. Write
$\rho=e^{2\pi\iu/3}$, so $\rho^2=\bar\rho$.

\subsection{Family 1, with an arithmetic twist}

The stuffle closed forms read
\begin{align}
  \Li_{a,b}(\rho,\rho)+\Li_{b,a}(\rho,\rho)     &=\Li_a(\rho)\Li_b(\rho)-\Li_{a+b}(\rho^2),
     \label{double:eq:eis-sym}\\
  \Li_{a,b}(\rho^2,\rho)+\Li_{b,a}(\rho,\rho^2) &=\Li_a(\rho^2)\Li_b(\rho)-\zeta(a+b),
     \label{double:eq:eis-mix}
\end{align}
with the diagonal corollary $\Li_{a,a}(\rho,\rho)=\tfrac12(\Li_a(\rho)^2-\Li_{2a}(\rho^2))$. The
arithmetic twist is in \eqref{double:eq:eis-mix}: because $\rho^2\cdot\rho=\rho^3=1$, the colored tail
degenerates to the \emph{ordinary} zeta value $\zeta(a+b)$, not a level-$3$ value. The single
cube-root $L$-values split cleanly under the distribution relation
$\Li_n(\rho)+\Li_n(\rho^2)=(3^{1-n}-1)\zeta(n)$:
\begin{equation}\label{double:eq:eis-single}
  \Real\Li_n(\rho)=\tfrac12(3^{1-n}-1)\,\zeta(n),\qquad
  \Imag\Li_n(\rho)=\Cl_n(2\pi/3),
\end{equation}
the real part a rational multiple of $\zeta(n)$ (or $\pi^n$ at even $n$), the imaginary part the
Clausen value---the level-$3$ analogue of $\beta(n)$, with $\Cl_2(2\pi/3)$ Gieseking's constant. All
of \eqref{double:eq:eis-sym}--\eqref{double:eq:eis-single} are confirmed in Wolfram~$15$.

\subsection{The $(\rho,\rho)$ frontier: one constant}

Running the level-$3$ double shuffle as in Section~\ref{double:sec:gauss-frontier}---$90$ stuffle and shuffle
relations from products $\Li_p(s)\Li_q(t)$, $s,t\in\{\rho,\rho^2\}$, each verified to
${\sim}10^{-600}$, exact rational rank---collapses the $(\rho,\rho)$ doubles still more sharply than
the Gaussian:

\begin{center}
\renewcommand{\arraystretch}{1.2}
\begin{tabular}{c c c c}
\hline
weight & unknowns (Re$+$Im) & dim (stuffle only) & dim (double shuffle)\\
\hline
$2$ & $2$  & $0$ & $0$\\
$3$ & $4$  & $2$ & $0$\\
$4$ & $6$  & $2$ & $0$\\
$5$ & $8$  & $4$ & $\mathbf 1$\\
$6$ & $10$ & $4$ & $0$\\
\hline
total & & $12$ & $\mathbf 1$\\
\hline
\end{tabular}
\end{center}

\begin{observation}\label{double:thm:eis-1}
Modulo single polylogarithms and the $(\rho,1)$ generators, the entire $(\rho,\rho)$ family through
weight~$6$ reduces to a \emph{single} candidate depth-two generator---$\Imag\Li_{3,2}(\rho,\rho)$, at
weight~$5$---with weights $2,3,4$ and $6$ collapsing completely.
\end{observation}

A gold-check at weight~$5$ fixes the survivor and predicts all seven other weight-$5$ components to
${\sim}10^{-696}$. The pattern mirrors the Gaussian: there too the weight-$5$ survivors are imaginary
$\Li_{3,2}$/$\Li_{4,1}$ directions.

\begin{remark}[a genuine cube-root obstacle]
The decisive shuffle $\Li_p(\rho^2)\Li_q(\rho)\to(\rho,\rho)+(\rho^2,\rho^2)$ lands only on points with
$w=z_1z_2\neq1$, so the computation never needs the mixed points $(\rho^2,\rho),(\rho,\rho^2)$, where
$w=\rho^3=1$. Those points presented a numerical obstacle in the earlier evaluator: both the
Lerch-plus-Wynn evaluator and a direct cumulative-sum evaluator stall there at ${\sim}8$ digits---the
period-three, non-oscillating tail defeats Wynn's $\varepsilon$-algorithm, unlike the even-period
$\iu,-1$ cases, which converge to full precision. Folding the mixed cube-root points in (for the
\emph{full} level-$3$ frontier) requires a period-three-aware accelerated evaluator, or Wolfram; the following mixed-point method addresses that computational gap.
\end{remark}
\section{Why the standard evaluator degenerates at \texorpdfstring{$w=1$}{w=1}}\label{mixed:sec:why}

With the convention the nested-series definition, the accelerated evaluator of the article writes the inner tail as
a Lerch transcendent and sums the \emph{outer} series with an acceleration factor $w^n=(z_1z_2)^n$. At
the mixed points $w=1$, that factor is identically $1$: the oscillation Wynn's $\varepsilon$-algorithm
relies on disappears, leaving a bare period-three power-law tail that the algorithm cannot resolve past
a handful of digits. Summing over the \emph{larger} index instead keeps an oscillating factor
$z_1^{n_1}$, but for $z_1=\rr^2$ this is a \emph{period-three} oscillation, and a plain
$\varepsilon$-algorithm (tuned for geometric tails) again stalls.

\section{A robust \texorpdfstring{$w=1$}{w=1} evaluator}\label{mixed:sec:eval}

Two ingredients fix it.

\paragraph{(a) A smooth correction identity.} Writing $l=n_1$ (the larger index) and splitting the
inner partial sum $P(l)=\sum_{n<l}z_2^n/n^b=\Li_b(z_2)-R(l)$ with tail
$R(l)=\sum_{n\ge l}z_2^n/n^b=z_2^{\,l}\,\Phi(z_2,b,l)$ (Lerch $\Phi$), one gets the exact identity
\begin{equation}\label{mixed:eq:smooth}
  \boxed{\;\Li_{a,b}(z_1,z_2)=\Li_a(z_1)\,\Li_b(z_2)\;-\;\sum_{l\ge1}\frac{w^{\,l}}{l^{a}}\,\Phi(z_2,b,l)\;}
\end{equation}
At $w=1$ the correction summand $\Phi(z_2,b,l)/l^a$ is \emph{smooth}---every oscillating factor
$z_1^l z_2^l=w^l$ collapses to $1$---with the clean asymptotic
$\Phi(z_2,b,l)\sim\sum_{j\ge0}\binom{-b}{j}c_j(z_2)\,l^{-b-j}$, where $c_0=\tfrac1{1-z_2}$ and
$c_{j\ge1}=\Li_{-j}(z_2)$ (rational in $z_2$). The subtle point---the source of a factor-of-one bug
that cost five digits in a first draft---is that $c_0=\sum_{k\ge0}z_2^k=\tfrac1{1-z_2}$ is \emph{not}
$\Li_0(z_2)=\tfrac{z_2}{1-z_2}$; they differ by exactly $1$. This gives an Euler--Maclaurin tail in
Hurwitz zetas, $\sum_{l>M}\Phi/l^a=\sum_j\binom{-b}{j}c_j\,\zeta(a+b+j,M+1)$, a fast and accurate
cross-check at low to medium precision.

\paragraph{(b) Period-aware $\rho$-acceleration for production.} For high precision the production
evaluator sums the outer cumulative series directly, samples the partial sums at the common period
$N=\lcm(\ord z_1,\ord z_2)$ (so every oscillating factor realigns to $1$ at the sample points), and
accelerates the resulting smooth, power-law-convergent sequence with \textbf{Wynn's $\rho$-algorithm}
---the algebraic-convergence counterpart of the $\varepsilon$-algorithm. The $\rho$-algorithm yields
empirical convergence on the tested families; precision inflation and independent comparison are needed to assess its conditioning
(\code{mpl\_w1.py}).

\section{Validation}\label{mixed:sec:valid}

The evaluator is checked by three independent, Wolfram-free routes plus one Wolfram cross-check.
\begin{enumerate}
\item \textbf{Exact symmetric closed form} (single polylogs only, an entirely separate code path):
  \begin{equation}\label{mixed:eq:sym}
    \Li_{a,b}(\rr^2,\rr)+\Li_{b,a}(\rr,\rr^2)=\Li_a(\rr^2)\Li_b(\rr)-\zeta(a+b),
  \end{equation}
  and the Gaussian analogue $\Li_{a,b}(\iu,-\iu)+\Li_{b,a}(-\iu,\iu)=\Li_a(\iu)\Li_b(-\iu)-\zeta(a+b)$.
  (The collision term is $\Li_{a+b}(z_1z_2)=\Li_{a+b}(1)=\zeta(a+b)$.)
\item \textbf{Conjugation} $\Li_{a,b}(\rr,\rr^2)=\overline{\Li_{a,b}(\rr^2,\rr)}$.
\item \textbf{Two independent summation methods}---the $\rho$-accelerated cumulative sum and the
  Euler--Maclaurin smooth-correction series---agree to full working precision.
\item \textbf{Wolfram~15} \texttt{MultiplePolyLog} at moderate precision.
\end{enumerate}
All four agree; the symmetric closed form \eqref{mixed:eq:sym} holds to the full generated precision
(${\sim}10^{-500}$) for every $(a,b)$.

\section{Closed forms for the Eisenstein mixed-point values}\label{mixed:sec:eisforms}

With $500$-digit values in hand, an integer-relation (PSLQ/LLL) search against the level-3 basis
---monomials in the single $L$-values $\{\pi,\log 3,\zeta(3),\zeta(5),\Cl_2(\tfrac{2\pi}3),
\Cl_4(\tfrac{2\pi}3),\Cl_6(\tfrac{2\pi}3)\}$ times the $(\rr,1)$ generators
$g_{a,b}=\Imag\Li_{a,b}(\rr,1)$, $h_{a,b}=\Real\Li_{a,b}(\rr,1)$, and the lone $(\rr,\rr)$
double-shuffle survivor $\Imag\Li_{3,2}(\rr,\rr)$---\textbf{reduces every mixed value through weight
$4$, and every weight-5 imaginary part}, to explicit closed forms (each canary-clean, height $\le 10$,
re-verified to $10^{-425}$ against the \emph{complete} weight-graded basis). The two simplest:
\begin{equation}\label{mixed:eq:eis-w2}
  \Imag\Li_{1,1}(\rr^2,\rr)=-\bigl(\Cl_2(\tfrac{2\pi}3)+\Imag\Li_{1,1}(\rr,1)\bigr),\qquad
  \Real\Li_{1,1}(\rr^2,\rr)=5\,\Real\Li_{1,1}(\rr,1)-\tfrac12\log^2 3.
\end{equation}
(The full list of ${\sim}16$ reductions is in \code{harvest\_eis\_mixed.py}.) The weight-5 \emph{real}
parts and all of weight $6$ do \emph{not} reduce to this basis at height $\lesssim 10^{4}$.

\section{The completed level-3 double-shuffle dimension}\label{mixed:sec:eisdim}

Two computations pin the dimension of the full level-3 depth-2 space---all four colored points
$(\rr,\rr),(\rr^2,\rr^2),(\rr^2,\rr),(\rr,\rr^2)$---modulo single polylogarithms and the $(\rr,1)$
generators.
\begin{itemize}
\item The \textbf{exact-rank convergent double shuffle} (stuffle and shuffle of products
  $\Li_p(s)\Li_q(t)$, $s,t\in\{\rr,\rr^2\}$, together with the convergent zeta-products
  $\Li_p(s)\zeta(q)$ whose shuffle reaches the mixed points) gives an \emph{upper bound} of $\mathbf 7$
  through weight $6$ (the divergent $\Li_1$ words, needing shuffle regularization, are omitted).
\item \textbf{Combining} that exact rank with the PSLQ closed forms of Section~\ref{mixed:sec:eisforms}
  collapses it:
\end{itemize}
\begin{center}
\begin{tabular}{ccl}
\toprule
weight & reported residual count & survivors\\
\midrule
$2$ & $0$ & ---\\
$3$ & $0$ & ---\\
$4$ & $0$ & ---\\
$5$ & $2$ & $\Imag\Li_{3,2}(\rr,\rr)$ \emph{(article)},\quad $\Real\Li_{4,1}(\rr^2,\rr)$ \textbf{(new)}\\
$6$ & $1$ & $\Imag\Li_{5,1}(\rr^2,\rr)$ \textbf{(new)}\\
\bottomrule
\end{tabular}
\end{center}
So the tested level-three quotient, after the numerical reductions, has reported residual count
$\mathbf 3$ through weight $6$. Everything through weight $4$ collapses entirely into the
$(\rr,1)$-generator world; the mixed cube-root points contribute \textbf{exactly two new constants}
beyond the article's single $(\rr,\rr)$ survivor---a \emph{real} direction $\Real\Li_{4,1}(\rr^2,\rr)$
at weight $5$ and an \emph{imaginary} direction $\Imag\Li_{5,1}(\rr^2,\rr)$ at weight $6$.

\begin{remark}[the article's own lesson]
Non-reduction by integer-relation search is strong evidence, not proof: a genuine reduction of large
height is invisible to PSLQ, and the convergent double shuffle omits the regularized $\Li_1$ relations.
So the two new survivors are best read as \emph{candidate} new level-3 depth-2 constants; regularized double shuffle can improve upper bounds; numerical independence would still need a separate theorem or a stated period conjecture. As an additional check, an \emph{enriched}-basis probe
(\code{probe\_eis\_candidates.py})---adding $\Li_n$ at candidate level-3 CM half-points
$\tfrac{1\pm\rr}2,\ \tfrac1{1-\rr},\ \tfrac{\rr}{\rr-1},\ 1-\rr$ (the level-3 analogue of the Gaussian
$\tfrac{1+\iu}2$ that resolved the article's Family-2)---fails to crack either candidate:
$\Real\Li_{4,1}(\rr^2,\rr)$ resists a $105$-atom basis (spurious scale ${\sim}10^{4}$).
\end{remark}

\section{The Gaussian \texorpdfstring{$(\iu,-\iu)$}{(i,-i)} point and an exact parallel}\label{mixed:sec:gauss}
\sloppy
The same evaluator handles the level-4 Gaussian point $(\iu,-\iu)$ that the article's level-4 run also
skipped ($w=\iu\cdot(-\iu)=1$). The picture is identical in shape.
\begin{itemize}
\item \textbf{Closed forms.} PSLQ against the level-4 basis (single $L$-values
  $\pi,\log 2,\beta(2),\beta(4),\beta(6),\zeta(3),\zeta(5)$, the half-point values
  $\Real\Li_n(\tfrac{1+\iu}2),\Imag\Li_n(\tfrac{1+\iu}2)$, and the $(\iu,1)$ generators) reduces every
  $(\iu,-\iu)$ value through weight $3$ and every weight-4 imaginary part. The weight-2 case is
  especially clean:
  \begin{equation}\label{mixed:eq:gauss-w2}
    \Li_{1,1}(\iu,-\iu)=-\Li_2\!\left(\tfrac{1-\iu}2\right)\qquad(\text{verified to } 10^{-50}).
  \end{equation}
\item \textbf{Dimension.} The corrected exact-rank double shuffle (now including $(\iu,-\iu)$) gives an
  upper bound of $12$; combined with the closed forms it collapses to the article's five $(\iu,-1)$
  constants plus \textbf{the same two new mixed directions}---$\Real\Li_{4,1}(\iu,-\iu)$ at weight $5$
  and $\Imag\Li_{5,1}(\iu,-\iu)$ at weight $6$. (One weight-3 residual, $\Imag\Li_{2,1}(\iu,-1)$, is a
  value the article already reduces to the half-point $\lambda/\mu$ basis, outside the convergent
  double shuffle's reach.)
\end{itemize}

\begin{observation}[the parallel]\label{mixed:thm:parallel}
In \emph{both} towers the $w=1$ mixed point leaves, beyond the earlier tested generators, two candidate mixed directions, and \emph{in the same two positions}:
\[
  \boxed{\ \Real\Li_{4,1}(\mathrm{mixed})\ \text{at weight }5,\qquad
            \Imag\Li_{5,1}(\mathrm{mixed})\ \text{at weight }6.\ }
\]
\end{observation}
That the level-3 and level-4 mixed points add new content at identical weights and index patterns
---a \emph{real} $\Li_{4,1}$ direction and an \emph{imaginary} $\Li_{5,1}$ direction---is the report's
most suggestive finding, and a natural target for the regularized double shuffle to explain.
